\documentclass[letterpaper]{article} % DO NOT CHANGE THIS
\usepackage[submission]{aaai2027}  % DO NOT CHANGE THIS
\usepackage[hyphens]{url}  % DO NOT CHANGE THIS
\usepackage{graphicx} % DO NOT CHANGE THIS
\urlstyle{rm} % DO NOT CHANGE THIS
\def\UrlFont{\rm}  % DO NOT CHANGE THIS
\usepackage{natbib}  % DO NOT CHANGE THIS AND DO NOT ADD ANY OPTIONS TO IT
\usepackage{caption} % DO NOT CHANGE THIS AND DO NOT ADD ANY OPTIONS TO IT
\frenchspacing  % DO NOT CHANGE THIS
\usepackage{amsmath}
\usepackage{amssymb}
\usepackage{amsthm}
\usepackage{algorithm}
\usepackage{algorithmic}
\usepackage{booktabs}
\usepackage{xspace}

\newtheorem{proposition}{Proposition}
\newtheorem{theorem}{Theorem}
\newtheorem{lemma}{Lemma}
\newtheorem{corollary}{Corollary}
\newtheorem{definition}{Definition}
\newtheorem{remark}{Remark}

% A result that also appears in the main text keeps the number it carries
% there, so the two documents can be read side by side. Results that exist only
% in this document use a separate S-series (Proposition S1, Theorem S1, ...),
% which cannot collide with a main-text number however either document grows.
\newtheorem{sproposition}{Proposition S}
\newtheorem{stheorem}{Theorem S}

\pdfinfo{
/TemplateVersion (2027.1)
}

\setcounter{secnumdepth}{2}

\title{Supplementary Material:\\EigenTree: a Centralizing Communication Protocol\\for Decentralized Multi-Agent Multi-Armed Bandit}
\author{
    Anonymous Submission
}
\affiliations{
    %
}

\begin{document}

\nocopyright

\maketitle

Notation follows the main paper throughout. $G=(V,E)$ is the underlying
graph, $N=|V|$, and $A$ its adjacency matrix. $\tilde\psi$ is the
(unnormalized, locally computed) leading eigenvector of $A$, with eigenvalue
$\lambda$, and $h^\star=\arg\max_i\tilde\psi_i$ is the hub. $\mathcal{T}$ is
the induced routing tree and $D$ its height. $\mathrm{par}(i)$, $\mathrm{ch}(i)$, and
$d_i$ denote a node's parent, children, and depth in $\mathcal{T}$, and
$\mathcal{N}_G(i)$ denotes node $i$'s neighbor set in $G$. We restate the one
standing assumption, since several results below are stated under it.

\begin{definition}[Expander regime]
\label{def:expander}
A family of graphs $\{G_N\}$ is in the expander regime if its algebraic
connectivity $\lambda_2$ (the second-smallest eigenvalue of the Laplacian) is
bounded below by a constant $c>0$ independent of $N$. We say a graph $G$ is in
the regime when it belongs to such a family, and treat $c$ as a constant in all
$N$-dependent bounds.
\end{definition}

\section{Hub Emergence, Routing, and Coordinator Sufficiency}

\begin{proposition}[Centralized aggregation and broadcast via routing]
\label{prop:aggbcast}
Suppose each node $i$ holds a scalar $x_i$ and propagates it toward
$h^\star$ one hop per round, with each intermediate node forwarding the sum
it received. Then the hub holds $\sum_i x_i$ after $D$ rounds, where $D$ is
the tree height. Symmetrically, a value broadcast from $h^\star$ one hop
per round reaches every node after $D$ rounds. Neither direction requires
any node to know $D$ in advance: the uplink terminates when the hub
observes a round with no incoming packets, and each node infers its own
depth $d_i$ from the hop count carried on the first downlink it receives.
\end{proposition}

\begin{proof}
\textbf{Uplink.} We show by induction on depth, from the leaves upward,
that every node at depth $d$ has forwarded the sum of its entire subtree's
values to its parent by the end of round $D-d+1$.

\emph{Base case.} A leaf $\ell$ at depth $d_\ell$ has no children, so its
subtree sum is simply $x_\ell$. It has nothing to wait for, so it forwards
$x_\ell$ to $\mathrm{par}(\ell)$ in round 1.

\emph{Inductive step.} Suppose every node at depth $d+1$ has forwarded its
subtree sum to its parent by the end of round $D-(d+1)+1 = D-d$. A node $i$
at depth $d$ with children $\mathrm{ch}(i)$ waits until it has received a
forwarded value from every child. Each child's value arrives no later than
round $D-d$ by the inductive hypothesis. Node $i$ sums these values
together with its own $x_i$ and forwards the total to $\mathrm{par}(i)$ in
round $D-d+1$.

Since $h^\star$ is at depth 0, the induction applied down to $d=1$ shows
every direct child of $h^\star$ has forwarded its subtree sum by round $D$.
The hub receives all of these by round $D$ and sums them with its own
$x_{h^\star}$ to obtain $\sum_i x_i$.

\emph{Termination without knowing $D$.} A node $i$ at depth $d$ cannot know
$D$, but it does not need to: it simply waits until every child, a fixed
and locally known set, has sent its value, then forwards. The hub likewise
does not know $D$ in advance. It accumulates incoming packets each round
and declares the uplink complete once a round passes with no new arrivals.
Every node forwards exactly once, by the induction above, so the hub
receives no further packets after round $D$, and this silence is a purely
local observation.

\textbf{Downlink.} A symmetric induction on depth, from the root downward,
gives the same result. $h^\star$ sends $v$ to $\mathrm{ch}(h^\star)$ in
round 1, tagged with hop count $0$. A node at depth $d$ receiving $(v,
\mathrm{hops}=d-1)$ sets its own depth $d_i \leftarrow (d-1)+1 = d$, which
is correct since it is exactly $d$ hops from $h^\star$ along tree edges,
and forwards $(v, \mathrm{hops}=d)$ to its own children in the next round.
By induction, every node at depth $d$ receives $v$ in round $d$, so every
node receives $v$ within $D$ rounds. Each node's depth is derived purely
from the hop count it personally receives, with no global knowledge
required.
\end{proof}

\begin{theorem}[Coordinator sufficiency]
\label{thm:coordinator}
Let $\mathcal{A}$ be any decision rule whose output at each step depends
only on globally aggregated per-agent statistics $B = \sum_{i \in V} b_i$,
where $b_i$ is agent $i$'s local contribution. Suppose a coordinator
receives the exact aggregate $B$ within $D$ rounds and broadcasts its
decision within $D$ further rounds, with all agents synchronized via the
depth signal carried on the broadcast. Then the coordinator's information
state at every decision point is identical to that of a centralized
instance of $\mathcal{A}$, with an additive overhead of at most $2D$ time
steps per decision round.
\end{theorem}

\begin{proof}
Fix any decision round $r$. By Proposition~\ref{prop:aggbcast} (uplink
direction, with $x_i = b_i$), after $D$ rounds the hub holds $B =
\sum_{i\in V} b_i$ exactly, not an approximation or a partial sum, since
every node's contribution is included by the termination argument above.
Because $\mathcal{A}$'s decision at round $r$ depends only on $B$, the hub
computes exactly the action a centralized instance of $\mathcal{A}$ would
compute given the same set of per-agent contributions.

By Proposition~\ref{prop:aggbcast} (broadcast direction), within a further
$D$ rounds every agent holds the hub's decision, and by the same mechanism
every agent has derived its own depth $d_i$ from the hop count on this
broadcast. The depth signal is carried on every downlink, so all agents can
identify the round at which the current cycle ends, namely $2D$ rounds
after the cycle's pull phase, and begin the next round's contribution
simultaneously, regardless of their individual depth in $\mathcal{T}$.

Across the whole run, the sequence of aggregates $B^{(1)}, B^{(2)},
\ldots$ seen by the hub is therefore identical, round for round, to the
sequence a centralized instance of $\mathcal{A}$ would see given the same
sequence of per-agent local contributions, so the sequence of decisions is
identical as well. The only difference from a literal centralized run is
the $2D$ rounds of communication latency inserted between successive
decisions: an additive overhead per round that does not alter the sequence
of decisions itself.
\end{proof}

\section{The MERW Motivation for Eigenvector Centrality}
\label{sec:merw}

The main text motivates electing by eigenvector centrality through a Maximal
Entropy Random Walk (MERW): the walk under which all paths of a given length
between two fixed nodes are equally likely. We give the standard construction
\citep{burda2009localization} and the identity the argument rests on, that
MERW's stationary distribution is proportional to the squared eigenvector
centrality, so its $\arg\max$ is the node EigenTree elects.

\begin{sproposition}[MERW stationary distribution]
\label{prop:merw}
Let $G$ be connected and non-bipartite, with adjacency matrix $A$, leading
eigenvalue $\lambda$, and positive unit leading eigenvector $\psi$. The
transition matrix
\[
P_{ij} \;=\; \frac{A_{ij}}{\lambda}\,\frac{\psi_j}{\psi_i}
\]
is stochastic, assigns equal probability to all paths of equal length between
any two fixed endpoints, and has stationary distribution $\pi_i = \psi_i^2$. In
particular $\arg\max_i \pi_i = \arg\max_i \psi_i = h^\star$.
\end{sproposition}

\begin{proof}
\emph{Stochastic.} Row $i$ sums to $\sum_j A_{ij}\psi_j/(\lambda\psi_i) =
(A\psi)_i/(\lambda\psi_i) = \lambda\psi_i/(\lambda\psi_i) = 1$, using
$A\psi=\lambda\psi$. Entries are nonnegative since $A\geq 0$ and $\psi>0$ by
Perron--Frobenius.

\emph{Equal path weights.} For a path $\gamma = (i_0, i_1, \ldots, i_L)$ with
every consecutive pair adjacent, the product telescopes:
\[
\prod_{t=0}^{L-1} P_{i_t i_{t+1}}
= \prod_{t=0}^{L-1} \frac{A_{i_t i_{t+1}}}{\lambda}\frac{\psi_{i_{t+1}}}{\psi_{i_t}}
= \frac{1}{\lambda^{L}}\,\frac{\psi_{i_L}}{\psi_{i_0}},
\]
since each $A_{i_t i_{t+1}}=1$ along an admissible path. The result depends only
on the endpoints $i_0, i_L$ and the length $L$, not on the route taken, so every
admissible path of length $L$ between the same two nodes carries the same
probability. This is the defining property the main text asks of the walk, and
it is what makes the ensemble non-local: an ordinary random walk weights each
step by $1/\deg$, giving longer weight to routes through low-degree nodes.

\emph{Stationarity.} With $\pi_i = \psi_i^2$, which is a distribution since
$\|\psi\|_2=1$,
\begin{align*}
\sum_i \pi_i P_{ij}
&= \sum_i \psi_i^2 \frac{A_{ij}}{\lambda}\frac{\psi_j}{\psi_i}
 = \frac{\psi_j}{\lambda}\sum_i A_{ij}\psi_i \\
&= \frac{\psi_j}{\lambda}\,\lambda\psi_j
 = \psi_j^2 = \pi_j,
\end{align*}
using the symmetry of $A$ and $A\psi=\lambda\psi$. Connected and non-bipartite
makes $P$ irreducible and aperiodic, so this stationary distribution is unique.

\emph{Argmax.} $t\mapsto t^2$ is strictly increasing on $t>0$ and $\psi>0$, so
squaring does not move the maximizer.
\end{proof}

The hub is thus the node carrying the largest share of the long-path traffic
under the ensemble in which no route is privileged, which is the sense in which
it is the graph's natural coordinator. Note the election needs only the
$\arg\max$, so the square is immaterial: agents compare $\tilde\psi_i$
directly and never form $\pi$.

\section{Cost of Not Using the Minimum Diameter Tree}
\label{sec:mdst-gap}

Throughout this section, $D = \mathrm{ecc}_G(h^\star)$ denotes the height
of the EigenTree routing tree, equivalently the eccentricity of $h^\star$ in
$G$. $D' = R(G) = \min_{v\in V}\mathrm{ecc}_G(v)$ denotes the radius of
$G$: the height achieved by rooting a shortest-path tree at $G$'s absolute
center, a minimum diameter spanning tree (MDST) \citep{hakimi1964,hassintamir1995}.

We first record that the elected hub is never a low-degree node, which is what
makes the pigeonhole hypothesis of Theorem~\ref{thm:mdst-gap} reachable: it
asks $\deg(h^\star)$ to be large, and electing by centrality already forces
$\deg(h^\star)\geq\lambda\geq\bar d$, the average degree.

\begin{sproposition}[Hub degree lower bound]
\label{prop:hubdeg}
The hub's degree is at least the graph's spectral radius: $\deg(h^\star) \geq \lambda$.
\end{sproposition}

\begin{proof}
$\tilde\psi$ is the (unnormalized) leading eigenvector of $A$ with
eigenvalue $\lambda$, so it satisfies $A\tilde\psi = \lambda\tilde\psi$.
Reading off the $h^\star$-th row of this equation,
\[
\lambda\tilde\psi_{h^\star} = \sum_{j \sim h^\star} \tilde\psi_j,
\]
where the sum ranges over the $\deg(h^\star)$ neighbors of $h^\star$ in
$G$. $\tilde\psi_{h^\star} = \max_k \tilde\psi_k$ by definition of
$h^\star$, so every term in the sum satisfies $\tilde\psi_j \leq
\tilde\psi_{h^\star}$, giving
\[
\lambda\tilde\psi_{h^\star} = \sum_{j\sim h^\star}\tilde\psi_j \leq \deg(h^\star)\cdot\tilde\psi_{h^\star}.
\]
By the Perron--Frobenius theorem, the leading eigenvector of a connected
graph's adjacency matrix is strictly positive, so $\tilde\psi_{h^\star} >
0$. We divide both sides by $\tilde\psi_{h^\star}$ to obtain $\lambda \leq
\deg(h^\star)$.
\end{proof}

\begin{stheorem}[Dense graphs: depth at most 2]
\label{thm:mdst-gap}
Let $\delta_{\bar N}(h^\star) = \min_{u \notin \mathcal{N}_G(h^\star)\cup\{h^\star\}} \deg(u)$ be the smallest degree among the nodes $h^\star$ is not directly connected to. If $\deg(h^\star) + \delta_{\bar N}(h^\star) \geq N$, then $D \leq 2$.
\end{stheorem}

\begin{proof}
Let $u \in V$ be any node with $u \neq h^\star$ and $u \notin
\mathcal{N}_G(h^\star)$: $u$ is not $h^\star$ itself and not a neighbor of
$h^\star$. We show $d(h^\star, u) \leq 2$. Since $u$ was an arbitrary
non-neighbor, and every neighbor of $h^\star$ trivially has
$d(h^\star,\cdot)=1$, this establishes $D = \mathrm{ecc}_G(h^\star) \leq 2$.

Consider $\mathcal{N}_G(h^\star)$ and $\mathcal{N}_G(u)$, the neighbor sets
of $h^\star$ and $u$ respectively. Both are subsets of $V \setminus
\{h^\star, u\}$. $h^\star \notin \mathcal{N}_G(h^\star)$ since there are no
self-loops, and $u \notin \mathcal{N}_G(h^\star)$ by assumption.
Symmetrically, $u \notin \mathcal{N}_G(u)$, and $h^\star \notin
\mathcal{N}_G(u)$ since $u \notin \mathcal{N}_G(h^\star)$ implies $h^\star
\notin \mathcal{N}_G(u)$ on an undirected graph. Hence
\begin{align*}
|\mathcal{N}_G(h^\star)| + |\mathcal{N}_G(u)| &= \deg(h^\star) + \deg(u) \\
&\geq \deg(h^\star) + \delta_{\bar N}(h^\star) \geq N,
\end{align*}
using $\deg(u) \geq \delta_{\bar N}(h^\star)$, since $u$ is by construction
one of the nodes over which $\delta_{\bar N}(h^\star)$ is minimized, and
the theorem's hypothesis.

$|V \setminus \{h^\star, u\}| = N - 2$, and both $\mathcal{N}_G(h^\star)$
and $\mathcal{N}_G(u)$ are subsets of this $(N-2)$-element set. Their sizes
sum to at least $N > N-2$, so they cannot be disjoint: this is a standard
pigeonhole argument, since $|X|+|Y| > |Z|$ for $X,Y\subseteq Z$ forces
$X\cap Y\neq\emptyset$, as otherwise $|X\cup Y| = |X|+|Y| > |Z|$ would
contradict $X\cup Y\subseteq Z$. So there exists $w \in
\mathcal{N}_G(h^\star)\cap\mathcal{N}_G(u)$, giving a length-2 path
$h^\star - w - u$, hence $d(h^\star,u)\leq 2$.
\end{proof}

Theorem~\ref{thm:mdst-gap} covers the dense case, where the hub reaches
everything in two hops. It says nothing about sparse expanders, which have
bounded degree, so no node reaches the whole graph in two hops and the
pigeonhole argument does not apply. There the gap is small for the opposite
reason: the MDST is itself deep, so a suboptimal root cannot lose much.

\begin{sproposition}[Sparse expanders: the MDST is logarithmically deep]
\label{prop:mdst-sparse}
If $G$ is in the expander regime of Definition~\ref{def:expander}, then $D' =
R(G) = \Theta(\log N)$, and consequently $D \leq 2D' = \Theta(\log N)$: even
the optimal root gives a logarithmically deep tree, and the elected hub is
within a factor $2$ of it.
\end{sproposition}

\begin{proof}
\emph{Upper bound.} A $\lambda_2$ bounded below by a constant forces
$\mathrm{diam}(G) = O(\log N)$ \citep{mohar1991}, and $R(G)\leq
\mathrm{diam}(G)$ always, so $D' = O(\log N)$.

\emph{Lower bound.} On any graph $\mathrm{diam}(G)\leq 2R(G)$: for any $u,w$
and a center $z$ achieving the radius, $d(u,w)\leq d(u,z)+d(z,w)\leq 2R(G)$ by
the triangle inequality. Hence $D'=R(G)\geq \mathrm{diam}(G)/2$. In the
expander regime the graph is connected with degrees bounded by $N-1$, and a
graph of maximum degree $\Delta$ has $\mathrm{diam}(G) = \Omega(\log_\Delta
N)$, since a ball of radius $r$ holds at most $1+\Delta\frac{(\Delta-1)^r -
1}{\Delta - 2}$ nodes and must cover all $N$. For bounded-degree expanders,
random regular graphs or ER just above the connectivity threshold, $\Delta =
O(1)$ or $O(\log N)$, giving $\mathrm{diam}(G) = \Omega(\log N /\log\log N)$
and so $D' = \Omega(\log N/\log\log N)$. Combined with the upper bound, $D'$ is
logarithmic up to the $\log\log$ factor, and exactly $\Theta(\log N)$ when
$\Delta = O(1)$.

\emph{Consequence.} $D = \mathrm{ecc}_G(h^\star) \leq \mathrm{diam}(G) \leq
2R(G) = 2D'$ for any root whatsoever, in particular $h^\star$. So the elected
hub never pays more than twice the optimal depth, and both are
$\Theta(\log N)$.
\end{proof}

The regime where the choice of root is genuinely expensive is the one
Definition~\ref{def:expander} excludes. On a path, $\lambda_2 = \Theta(N^{-2})$
and $\mathrm{diam}(G) = \Theta(N)$, so the factor-$2$ bound above is still
true but both quantities are linear in $N$.

\section{Instantiations: Regret and Best Arm Identification}

\begin{corollary}[Regret matches centralized UCB]
\label{cor:regret}
Under the EigenTree relay protocol, the leading term of the expected group
cumulative regret matches a centralized UCB agent on $P$ total pulls, with
$P = \Theta(NT/(2D+1))$.
\end{corollary}

\begin{proof}
Instantiate Theorem~\ref{thm:coordinator} with $b_i$ as agent $i$'s new
per-arm pull counts and reward sums since the last cycle. The hub's
aggregate state $(\hat n^k, \hat s^k)$ for each arm $k$, together with the
total-pull count $P$, is then identical at every decision point to the
state a centralized UCB agent would hold after observing the same pulls.
$P$ itself is maintained exactly: the hub increments it by the number of
uplink contributions received plus its own pull each cycle, which by
Proposition~\ref{prop:aggbcast} is exactly the number of new pulls made
network-wide since the last decision. The standard UCB regret bound is a
function only of $(\hat n^k,\hat s^k,P)$ evaluated at the algorithm's own
internal clock $P$, not wall-clock time, and this sequence of internal
states is identical between the relay protocol and a literal centralized
run. So the leading-order regret term is identical as a function of $P$.

The relationship between $P$ and wall-clock horizon $T$ follows from the
fixed cycle structure. Each cycle lasts $2D+1$ rounds: one simultaneous
pull round, $D$ uplink rounds, one hub-decision round, $D$ downlink
rounds. Every cycle contributes exactly $N$ new pulls, one per agent, to
$P$. Over horizon $T$, the number of completed cycles is
$\Theta(T/(2D+1))$, so $P = \Theta(NT/(2D+1))$.
\end{proof}

\begin{corollary}[BAI sample complexity matches centralized BAI]
\label{cor:bai}
The total number of arm pulls required by EigenTree-BAI to identify $a^\star$ with probability $\geq 1-\delta$ matches the centralized algorithm of \citet{hillel2013distributed}; the only overhead is $2D+1$ extra time steps per elimination round.
\end{corollary}

\begin{proof}
EigenTree-BAI uses exactly Hillel's Algorithm 3 pull schedule $L_r = \lceil t_r
- t_{r-1}\rceil$ and elimination rule, eliminate arms with estimated mean
below the round-$r$ threshold $\epsilon_r$, applied to the hub's aggregate
statistics rather than an all-to-all broadcast. Instantiate
Theorem~\ref{thm:coordinator} with $b_i$ as agent $i$'s per-arm sums over
the current surviving arm set $S_{r-1}$. The hub's aggregate after each
round's uplink is then exactly the sum of pulls Hillel's centralized
algorithm would have observed given the same schedule and elimination
history, since every agent executes the identical, schedule-determined
number of pulls $L_r$ per round. The schedule depends only on
$N,K,\delta,r$, all known in advance, so no agent needs to learn anything
to determine its own pull count. The elimination decision depends only on
this aggregate, so the sequence of surviving sets $S_1 \supseteq S_2
\supseteq \cdots$ produced by EigenTree-BAI is identical to Hillel's, round for
round, and the total number of pulls until $|S_r|=1$ is identical.

The only addition is communication. Each round $r$ requires $D$ uplink
rounds to deliver the round's local sums to the hub, one round for the hub
to apply the elimination rule, and $D$ downlink rounds to broadcast the
surviving set $S_r$ and the depth signal: $2D+1$ rounds of overhead per
elimination round beyond the $L_r$ pull rounds themselves.
\end{proof}

\section{Fault Tolerance}
\label{sec:ft}

Throughout this section, $\mathcal{T}$, $D$, $\mathrm{par}(\cdot)$,
$\mathrm{ch}(\cdot)$, $d_i$ refer to the routing tree and its structure
\emph{before} the failure under discussion. A subscript or prime denotes
the corresponding quantity afterward, as introduced at each result.

Re-election needs the graph to survive the failure. This is supplied by the
expander regime of Definition~\ref{def:expander}, not assumed separately:
$\kappa(G)\geq\lambda_2$ \citep{fiedler1973} bounds the vertex connectivity
below by the algebraic connectivity, so a $\lambda_2$ bounded away from $0$ by a
constant $c>0$ gives $\kappa(G)\geq c$, and for $N$ large enough that $c>1$ the
graph stays connected after any single node is removed. The one condition of
Definition~\ref{def:expander} therefore makes re-election both possible, since
$G\setminus\{v\}$ is connected, and cheap, since the same $\lambda_2$ sets the
contraction rate.

\subsection{Re-Election}

Deleting node $v$ from $A$ gives $\tilde A = A + E$ with
$E = -(A e_v e_v^\top + e_v e_v^\top A)$. As $A$ has no self-loops, $E$ has rank
at most $2$ and $\|E\|_2 = \|A e_v\|_2 = \sqrt{\deg(v)}$, since $A e_v$ is the
$v$-th column of $A$, a $0/1$ vector with $\deg(v)$ ones. This holds for any $v$;
the hub is only the case of largest $\deg(v)$.

\begin{proposition}[Warm-started re-election]
\label{prop:reelection}
After any single node failure, the warm-started power iteration on $\tilde A$ converges to the leading eigenvector of $\tilde A$ from the pre-failure $\tilde\psi$. After the fixed budget $\tau_{\mathrm{init}} = O(\log N)$, the residual angle is at most $(\lambda_2/\lambda_1)^{\tau_{\mathrm{init}}}\cdot\sqrt{\deg(v)}/\gamma$, where $\gamma$ is the eigenvalue gap of $\tilde A$, against $(\lambda_2/\lambda_1)^{\tau_{\mathrm{init}}}\cdot\Theta(1)$ for a cold start from uniform. The warm start is therefore never worse, and its advantage is a factor $\gamma/\sqrt{\deg(v)}$, largest when the failed node has low degree.
\end{proposition}

\begin{proof}
After $v$ goes silent, a surviving node's update
$\tilde w_i = \sum_{j\sim i} w_j$ no longer includes the dead neighbor, so the
operator applied each round is exactly $\tilde A$ on the live nodes. Since $G$ is
in the expander regime, $G\setminus\{v\}$ is connected by the argument above, so
$\tilde A$ is irreducible and has
a unique positive leading eigenvector $\tilde\psi_{\tilde A}$ by
Perron--Frobenius, and power iteration converges to it from any start with
nonzero overlap; the warm start $\tilde\psi$ qualifies, as both vectors are
strictly positive.

For the accuracy, the initial angle to the target is bounded by Davis--Kahan
\citep{daviskahan1970},
$\sin\Theta(\tilde\psi,\tilde\psi_{\tilde A}) \leq \|E\|_2/\gamma
= \sqrt{\deg(v)}/\gamma$, using $\|E\|_2 = \sqrt{\deg(v)}$ for the rank-2
perturbation $E$ deleting row and column $v$. Each power-iteration round
multiplies the component orthogonal to $\tilde\psi_{\tilde A}$ by at most
$\lambda_2/\lambda_1$ relative to the leading one, so the residual angle after
$t$ rounds is at most $(\lambda_2/\lambda_1)^t\tan\Theta$.

The survivors run the fixed budget $t = \tau_{\mathrm{init}} = O(\log N)$
rather than a tolerance-driven number of rounds. This is forced: a budget
derived from the bound above would require $\deg(v)$, which only $v$'s
neighbors know and which $v$ cannot report once dead, and $\gamma$, a spectral
quantity of $\tilde A$ that no node holds. Every node can compute
$\tau_{\mathrm{init}}$ from $N$ alone, so all survivors finish on the same
round, exactly as at initialization.

Substituting $t = \tau_{\mathrm{init}}$ and the two initial angles gives the
claim: the warm start begins at $\tan\Theta \leq \sqrt{\deg(v)}/\gamma$, while
a cold start from the uniform vector begins at $\tan\Theta = \Theta(1)$, having
no dependence on the failure. Both contract at the same rate over the same
budget, so the warm start's residual is smaller by the factor
$\gamma/\sqrt{\deg(v)}$, and it is never worse. The advantage is largest when
$\deg(v)$ is small; when $\deg(v) = \Theta(N)$ it narrows to
$\gamma/\sqrt{N}$, since $\deg(v) \leq N$ always.
\end{proof}

\begin{remark}
\label{rem:gamma}
The bound treats $\gamma$, the eigenvalue gap of the \emph{perturbed} matrix
$\tilde A$, as $\Omega(1)$ in the expander regime. This is a property of
$\tilde A$, not of $A$, so it requires the gap to survive the failure. Since
$E$ has rank at most $2$, Weyl's inequality moves each eigenvalue of $A$ by at
most $\|E\|_2 = \sqrt{\deg(v)}$, which bounds the perturbed gap in terms of the
original one but does not by itself keep it away from $0$ when $\deg(v)$ is
large. A failure-independent lower bound on $\gamma$ is left for future work.
\end{remark}

\subsection{Detection and Synchronous Restart}
\label{sec:ft-detect}

Before the survivors can re-elect they must detect the failure and enter the new
power iteration in step, with no coordinator and no synchronized clocks. Both
follow from the routing tree already in place. Every node learned its depth $d_i$
and the height $D$ from the first downlink, where depth is the hop count that
message carried, and every ordinary cycle is one uplink-and-downlink traversal of
length $2D+1$ hops. Each node therefore knows, for each of its tree neighbors, the
exact hop of the cycle at which that neighbor's message is due, and runs a single
uniform timer $\tau = 2D+1$ hops against that schedule, suspecting a neighbor dead
when its due message fails to arrive within $\tau$. The rule is identical for
every node, and no node is a distinguished detector: the hub is detected by its
children exactly as any node is detected by its parent.

\begin{proposition}[Failure horizon]
\label{prop:horizon}
Under the uniform timer $\tau = 2D+1$, after any single node failure every surviving node receives the failure notice within $2D$ hops of the earliest detection.
\end{proposition}

\begin{proof}
Removing $v$ severs $\mathcal{T}$ into components: the component above $v$ (still
containing every node not below $v$) and one component per child
$c\in\mathrm{ch}(v)$, since every tree path between two of these components ran
through $v$. No survivor can flood all of them over tree edges, because the paths
between components are exactly the ones that passed through $v$; instead each
component has its own first detector, the parent $\mathrm{par}(v)$ for the upper
component and the child $c$ for each lower one, and each detector floods the
notice over the intact tree edges of its own component only.

Two quantities bound the horizon. First, the detectors do not fire
simultaneously: a detector at depth $d$ expects the failed node's message at hop
$d$ (uplink) or hop $2D-d$ (downlink) of the cycle, so two detectors' firing
times differ by their difference in these expected hops, which is at most the
depth spread $D$. Since all detectors use the same $\tau$ measured against the
same round-1 schedule, the last detector fires at most $D$ hops after the first.
Second, within a component the notice floods over tree edges from the detector,
reaching the farthest node in at most the component's height, itself at most $D$
hops. Adding the two, the last node to hear does so at most $D + D = 2D$ hops
after the earliest detection.
\end{proof}

The same $2D$ budget makes the restart synchronous. The notice is flooded
carrying the downlink's hop counter, so a node that receives it $h$ hops from its
detector knows that $h$ of the $2D$ hops have elapsed and waits the remaining
$2D-h$ hops before taking its first power-iteration step. Every node's
$(\text{hops travelled}) + (\text{hops waited}) = 2D$, so all survivors take
their first step at the same instant, measured from the earliest detection, with
no clock agreement. The wait is read directly off the hop counter, which is the
node's depth $d_i$ up to the fixed offset it always was, so the depth each node
already holds serves three roles at once, the detection deadline through $\tau$,
the bound on detector skew, and the restart countdown, and no field is added to
any message.

\subsection{Cost of a Recovery}
\label{sec:recovery-cost}

The main text reports recovery times measured from the failure injection to
the round at which ordinary cycling resumes. That total decomposes into four
stages, each already established above, and the decomposition is what those
measurements are checked against.

\begin{sproposition}[Recovery cost]
\label{prop:recovery-cost}
After a single node failure, the survivors resume ordinary cycling within
\[
\underbrace{(2D+1)}_{\text{detection}} \;+\; \underbrace{2D}_{\text{notice}}
\;+\; \underbrace{\tau_{\mathrm{init}}}_{\text{re-election}}
\;+\; \underbrace{2D_{\mathrm{new}}}_{\text{rebuild}}
\]
rounds of the failure, where $D$ is the height of the pre-failure tree,
$D_{\mathrm{new}}$ that of the rebuilt one, and $\tau_{\mathrm{init}} =
O(\log N)$.
\end{sproposition}

\begin{proof}
The four stages run in sequence, so their costs add.

\emph{Detection.} Each node times its tree neighbors against the cycle
schedule with the uniform timer $\tau = 2D+1$ of Section~\ref{sec:ft-detect}.
A node whose expected message misses that window is suspected dead, so the
first detector fires at most $2D+1$ rounds after the failure.

\emph{Notice.} By Proposition~\ref{prop:horizon}, every survivor holds the
notice within $2D$ hops of the earliest detection, and the hop counter carried
on the notice makes them all take their first iteration step on the same round.

\emph{Re-election.} The survivors run the warm-started power iteration of
Proposition~\ref{prop:reelection} for the fixed budget $\tau_{\mathrm{init}} =
O(\log N)$, computed from $N$ alone, so they finish together.

\emph{Rebuild.} Max-flooding over $G\setminus\{v\}$ re-roots the tree at the
new argmax exactly as at initialization. The first ordinary cycle is one
uplink-and-downlink traversal of the rebuilt tree, so $2D_{\mathrm{new}}$
rounds pass before the new height is known and cycling is back in step.

Summing the four gives the claim. Only $\tau_{\mathrm{init}}$ is spectral;
the other three scale with tree height, which is why a recovery costs more
when the hub moves to a node of larger eccentricity.
\end{proof}

\section{Additional Experiments}
\label{sec:exp}

The main text evaluates EigenTree on Erd\H{o}s--R\'{e}nyi graphs, the setting
its analysis targets: above the connectivity threshold an ER graph is an
expander with high probability, so it lies in the regime of
Definition~\ref{def:expander} where the costs are stated. Nothing in the
protocol is specific to that model, so this section repeats both
instantiations on three further topologies: Barab\'{a}si--Albert, a $2$D grid,
and a barbell. Together they vary degree distribution, diameter, and the
sharpness of the centrality peak.

\subsection{The Elected Hub and Routing Tree}
\label{sec:exp-trees}

Figures~\ref{fig:tree-er}--\ref{fig:tree-barbell} show what the election
produces on each of the four graphs. The left panel colours every node by its
eigenvector centrality $\psi_i/\psi_{\max}$, the middle panel overlays the
routing tree on the graph, and the right panel redraws that tree as a
hierarchy with every arrow pointing from a child to its parent, so the whole
structure funnels into the hub.

The four cases differ in how sharply the centrality peaks. On the
Barab\'{a}si--Albert graph one high-degree node dominates and the tree is
shallow and wide, most nodes attaching to the hub directly. The grid is
near-regular, so the centrality is spread almost evenly, the hub sits at the
centre of the lattice, and the tree is correspondingly deep. The barbell
splits the mass between the two clique-bridge endpoints; the hub is one of
them and reaches the far clique through the other, which is why its depth is
$2$ rather than the $3$ its diameter would suggest. The
Erd\H{o}s--R\'{e}nyi graph sits between these, with a mild peak and a
two-hop tree.

\begin{figure*}[t]
    \centering
    \includegraphics[width=\textwidth]{results/MERW_visualization/merw_tree_er_N20_p0.5.pdf}
    \caption{Centrality and routing tree on an Erd\H{o}s--R\'{e}nyi graph
    ($N=20$, $p=0.5$).}
    \label{fig:tree-er}
\end{figure*}

\begin{figure*}[t]
    \centering
    \includegraphics[width=\textwidth]{results/MERW_visualization/merw_tree_ba_N20.pdf}
    \caption{Centrality and routing tree on a Barab\'{a}si--Albert graph
    ($N=20$, $m=2$). The centrality concentrates on one high-degree node, so
    the tree is shallow and wide.}
    \label{fig:tree-ba}
\end{figure*}

\begin{figure*}[t]
    \centering
    \includegraphics[width=\textwidth]{results/MERW_visualization/merw_tree_grid_N25.pdf}
    \caption{Centrality and routing tree on a $5\times5$ grid. The graph is
    near-regular, so the centrality is nearly flat and the elected hub sits at
    the centre of the lattice.}
    \label{fig:tree-grid}
\end{figure*}

\begin{figure*}[t]
    \centering
    \includegraphics[width=\textwidth]{results/MERW_visualization/merw_tree_barbell_N20.pdf}
    \caption{Centrality and routing tree on a barbell graph ($N=20$). The hub
    is one of the two bridge endpoints and reaches the far clique through the
    other, giving depth $2$ on a graph of diameter $3$.}
    \label{fig:tree-barbell}
\end{figure*}

\subsection{Best Arm Identification Across Topologies}
\label{sec:exp-bai}

Corollary~\ref{cor:bai} predicts that EigenTree-BAI needs exactly as many arm
pulls as the centralized algorithm of \citet{hillel2013distributed} on any
graph where the tree spans the network: the hub's aggregate is exact by
Proposition~\ref{prop:aggbcast}, and the elimination schedule $L_r$ depends
only on $N,K,\delta,r$. The topology enters through $D$ alone, which is a cost
in time steps, not in samples.

Figure~\ref{fig:bai-topologies} bears this out. The two algorithms draw
independent reward streams, so the comparison is an ordinary two-sample one;
their means differ by $8.9\%$ on average and at most $0.29$ standard errors at
any size, well inside the run-to-run spread. A fixed-confidence stopping time
is itself highly variable, with a coefficient of variation near $0.7$ here, so
differences of this size are noise rather than a cost of routing. The first
three panels coincide because at fixed $N$ the pull count is a function of the
elimination schedule and the rewards drawn, and neither knows the graph; the
grid differs only because it is sized by its side length and so realizes
$\lfloor\sqrt N\rfloor^2$ agents rather than $N$.

\begin{figure*}[t]
    \centering
    \includegraphics[width=\textwidth]{results/supplementary/bai_topologies.pdf}
    \caption{Total arm pulls to identify the best arm, EigenTree-BAI against
    the centralized reference \citep{hillel2013distributed}, on four
    topologies ($K=10$, $\delta=0.05$, 50 runs per $N$). The two draw
    independent reward streams; error bars are the standard deviation across
    runs. Their means differ by $8.9\%$ on average and at most $0.29$ standard
    errors at any size, so the curves are indistinguishable given how variable
    a fixed-confidence stopping time is. The first three panels coincide
    because at fixed $N$ the pull count is a function of the elimination
    schedule and the rewards drawn, and neither knows the graph; the grid
    differs only because it realizes $\lfloor\sqrt N\rfloor^2$ agents. This is
    what Corollary~\ref{cor:bai} predicts: the tree costs communication,
    $2D+1$ rounds per elimination round, rather than samples.}
    \label{fig:bai-topologies}
\end{figure*}

\subsection{Regret Minimization Across Topologies}
\label{sec:exp-regret}

Corollary~\ref{cor:regret} ties EigenTree-UCB's regret to a centralized UCB
agent run on $P$ total pulls, with $P = \Theta(NT/(2D+1))$. The topology again
enters only through $D$, so the prediction is that regret tracks the
centralized baseline as a function of pulls, while the number of pulls fitting
into a fixed horizon falls as the tree deepens.

Figure~\ref{fig:regret-topologies} therefore plots regret against total pulls
rather than against the round. The distinction matters: a cycle costs $2D+1$
rounds, so a deeper tree completes fewer cycles and makes fewer pulls in the
same $T$, and its cumulative regret at $T$ would read lower merely because it
has played less. Drawing every curve out to the largest budget all three
algorithms reach removes that confound.

On that footing the result is uniform. The EigenTree-UCB and Central-UCB
curves lie on top of each other on every topology, within bands that overlap
throughout, so the two are indistinguishable at this sample size rather than
one edging the other. Coop-UCB2 sits clearly above both everywhere.

Its grid curve stands out by two orders of magnitude, which is why that panel
is drawn on a logarithmic ordinate. A small spectral gap alone does not
account for it, since the barbell's $\lambda_2$ is smaller still ($0.039$
against the grid's $0.098$ at $N=100$) and its Coop-UCB2 regret is an order of
magnitude lower; what separates the grid is that low degree and large diameter
apply together, so the doubly stochastic weights both move little mass per
round and must move it far. EigenTree is insensitive to either, since
Proposition~\ref{prop:aggbcast} delivers the exact aggregate in $D$ hops
however slowly the graph mixes, and that insensitivity is the practical case
for routing over a tree instead of averaging over the graph. The price
EigenTree pays is visible in how far each panel extends along the pull axis,
which shrinks as $D$ grows.

\begin{figure*}[t]
    \centering
    \includegraphics[width=\textwidth]{results/supplementary/regret_topologies.pdf}
    \caption{Group cumulative regret against total arm pulls on four
    topologies ($N=100$, $K=5$, $T=5000$, 50 runs); bands are the standard
    deviation across runs. Plotting against pulls rather than rounds compares
    the algorithms on equal samples, which keeps the depth cost in time where
    Corollary~\ref{cor:regret} puts it: a cycle costs $2D+1$ rounds, so a
    deeper tree simply reaches fewer pulls within the horizon. EigenTree-UCB
    tracks Central-UCB on every topology, while Coop-UCB2 is worse throughout.
    The grid panel uses a logarithmic ordinate, since Coop-UCB2 runs two
    orders of magnitude above the other two there.}
    \label{fig:regret-topologies}
\end{figure*}

\section{Reproducibility}
\label{sec:reproducibility}

\paragraph{Computing infrastructure.} All experiments run on CPU only, on an
Apple M1 Pro with 16\,GB of RAM under macOS~26.5. The code is pure Python and
uses no GPU.

\paragraph{Software.} Python~3.12 with NumPy~2.4, NetworkX~3.6, and
Matplotlib~3.10; no other third-party libraries are required. The three scripts
in the code archive (\texttt{regret\_min.py}, \texttt{best\_arm\_id.py},
\texttt{fault\_tolerance.py}) are self-contained and share no imports. Each
takes a \texttt{--graph} argument selecting the topology, so every table in
Section~\ref{sec:exp} is reproduced by rerunning the corresponding script once
per topology.

\paragraph{Randomness and seeds.} Every experiment is fully deterministic given
its seed. A master \texttt{RandomState} seeded at a fixed value derives one graph
seed, held constant across all runs so every algorithm is compared on the same
graph, and one independent per-run seed; each run reseeds the generator from its
per-run seed before drawing rewards. Graph generation, arm means, and reward
draws all flow from these seeds, so rerunning a script reproduces the reported
curves exactly.

Gossip partner selection during the power iteration draws from a generator of
its own, separate from the one supplying rewards. This matters for the
like-for-like comparisons: electing the hub consumes randomness, so drawing
gossip partners from the reward stream would leave the protocol and its
baseline reading different rewards from the same seed, and they would then
diverge through an artifact of the simulation rather than through anything the
protocol does.

In the BAI comparison each algorithm additionally draws its own reward stream.
Sharing one would couple them exactly: both request the same arms in the same
order, so they would receive identical rewards, eliminate identically, and
report the same pull count with zero variance in the difference. That is a
valid check that routing loses no information, but it is not a measurement of
sample complexity, since the agreement would hold by construction rather than
be observed. Independent streams make Figure~\ref{fig:bai-topologies} an
ordinary two-sample comparison.

\paragraph{Parameters.} The main-text experiments use the settings stated
there: for regret minimization, $N=100$, $K=5$, $T=5000$, $p=0.5$, 50 runs; for
best arm identification, $K=10$, $\delta=0.05$, $p=0.5$, 50 runs per $N \in
\{10,20,30,40,50\}$; for fault tolerance, $N=20$, $p=0.5$, $K=10$, $T=2000$,
with the hub killed at $t=500,1000,1500$. The UCB exploration constant is $c=2$
and the sub-Gaussian parameter is $\sigma=1$. Reported quantities are means over
the runs, with standard error of the mean shown as shaded bands or error bars.

The topology sweeps of Section~\ref{sec:exp} reuse these settings and vary only
the graph. Grid graphs are built with side $\lfloor\sqrt N\rfloor$, so their
realized agent count is read off the generated graph rather than the requested
$N$; both the protocol and its baseline are given that same count, so the
comparison stays like-for-like.

\bibliography{references}

\end{document}
